\subsection{The Semantic Web}
\label{sec:semantic-web}
\slidesrc{Ontologies: 8, 39--48}

\subsubsection{The vision}
\slidesrc{Ontologies: 8, 40--41}

\begin{nbdef}{Semantic Web}
``The \textbf{Semantic Web} is not a separate Web but an extension of the current one,
in which information is given a well-defined meaning, better enabling computers and
people to work in cooperation.'' (Tim Berners-Lee, James Hendler and Ora Lassila,
\emph{Scientific American}, May 2001.)
\end{nbdef}

The same article continues: the Semantic Web will bring structure to the meaningful
content of Web pages, creating an environment where \textbf{software agents} roaming
from page to page can readily carry out sophisticated tasks for users. The goals are
\begin{itemize}[nosep]
  \item automatic \textbf{discovery} of Web resources and Web services;
  \item automatic \textbf{composition} of Web resources and Web services;
  \item automatic \textbf{execution and monitoring} of Web services.
\end{itemize}

\begin{nbintuition}[why the Web needs ontologies]
A page saying ``Dr.\ Rossi, Tuesday 10--12, room 506'' is perfectly clear to a human
and opaque to a program: is it an office hour, a lecture, an appointment? If the page
also states, in RDF, that this is an instance of \emph{OfficeHour} with properties
\emph{teacher}, \emph{day} and \emph{room}, defined in a shared ontology, an agent can
find all office hours of a department, compare them with your calendar and book one.
Ontologies give the shared vocabulary; RDF gives the format of the facts.
\end{nbintuition}

\subsubsection{The Semantic Web stack}
\slidesrc{Ontologies: 42--45}

The Semantic Web technologies are usually drawn as a \textbf{stack of layers}, each
built on the ones below, like the language evolution of
Figure~\ref{fig:languages-evolution}. The stack itself evolved over the years.

\begin{center}
\begin{tikzpicture}[
    x=1.3cm, y=0.55cm]
  % ---- 2003 stack
  \fill[nbgray!25] (0,0) rectangle (1.95,1);
  \node[font=\scriptsize] at ($(0,0)!0.5!(1.95,1)$) {URI};
  \fill[nbgray!25] (2.05,0) rectangle (4,1);
  \node[font=\scriptsize] at ($(2.05,0)!0.5!(4,1)$) {Unicode};
  \fill[nborange!25] (0,1.1) rectangle (1.95,2.1);
  \node[font=\scriptsize] at ($(0,1.1)!0.5!(1.95,2.1)$) {XML};
  \fill[nborange!25] (2.05,1.1) rectangle (4,2.1);
  \node[font=\scriptsize] at ($(2.05,1.1)!0.5!(4,2.1)$) {Namespaces};
  \fill[nborange!45] (0.4,2.2) rectangle (3,3.2);
  \node[font=\scriptsize] at ($(0.4,2.2)!0.5!(3,3.2)$) {RDF M\&S};
  \fill[nborange!45] (0.6,3.3) rectangle (3,4.3);
  \node[font=\scriptsize] at ($(0.6,3.3)!0.5!(3,4.3)$) {RDF Schema};
  \fill[nbblue!20] (0.8,4.4) rectangle (3,5.4);
  \node[font=\scriptsize] at ($(0.8,4.4)!0.5!(3,5.4)$) {Ontology};
  \fill[nbgreen!25] (1.2,5.5) rectangle (3,6.5);
  \node[font=\scriptsize] at ($(1.2,5.5)!0.5!(3,6.5)$) {Rules};
  \fill[nbgreen!25] (1.0,6.6) rectangle (3,7.8);
  \node[font=\scriptsize] at ($(1.0,6.6)!0.5!(3,7.8)$) {Logic framework};
  \fill[nbteal!25] (1.8,7.9) rectangle (3,8.9);
  \node[font=\scriptsize] at ($(1.8,7.9)!0.5!(3,8.9)$) {Proof};
  \fill[nbpurple!20] (2.0,9.0) rectangle (3.9,10);
  \node[font=\scriptsize] at ($(2.0,9.0)!0.5!(3.9,10)$) {Trust};
  \fill[nbred!25] (3.1,2.2) rectangle (3.5,8.9);
  \node[font=\scriptsize, rotate=90] at ($(3.1,2.2)!0.5!(3.5,8.9)$) {Signature};
  \fill[nbred!25] (3.55,2.2) rectangle (3.95,8.9);
  \node[font=\scriptsize, rotate=90] at ($(3.55,2.2)!0.5!(3.95,8.9)$) {Encryption};
  \node[font=\small\bfseries] at (2,-0.8) {2003};
  % ---- 2005 stack
  \begin{scope}[xshift=7cm]
  \fill[nbgray!25] (0,0) rectangle (1.95,1);
  \node[font=\scriptsize] at ($(0,0)!0.5!(1.95,1)$) {URI};
  \fill[nbgray!25] (2.05,0) rectangle (4,1);
  \node[font=\scriptsize] at ($(2.05,0)!0.5!(4,1)$) {Unicode};
  \fill[nborange!25] (0,1.1) rectangle (1.95,2.1);
  \node[font=\scriptsize] at ($(0,1.1)!0.5!(1.95,2.1)$) {XML};
  \fill[nborange!25] (2.05,1.1) rectangle (4,2.1);
  \node[font=\scriptsize] at ($(2.05,1.1)!0.5!(4,2.1)$) {Namespaces};
  \fill[nbred!45] (0,2.2) rectangle (3,3.2);
  \node[font=\scriptsize] at ($(0,2.2)!0.5!(3,3.2)$) {RDF Core};
  \fill[nborange!45] (0.6,3.3) rectangle (3,4.3);
  \node[font=\scriptsize] at ($(0.6,3.3)!0.5!(3,4.3)$) {RDF Schema};
  \fill[nborange!45] (0.6,4.4) rectangle (3,5.4);
  \node[font=\scriptsize] at ($(0.6,4.4)!0.5!(3,5.4)$) {DLP bit of OWL/Rules};
  \fill[nbblue!20] (0.6,5.5) rectangle (1.75,6.5);
  \node[font=\scriptsize] at ($(0.6,5.5)!0.5!(1.75,6.5)$) {OWL};
  \fill[nbgreen!25] (1.85,5.5) rectangle (3,6.5);
  \node[font=\scriptsize] at ($(1.85,5.5)!0.5!(3,6.5)$) {Rules};
  \fill[nbblue!35] (0,3.3) rectangle (0.5,6.5);
  \node[font=\scriptsize, rotate=90] at ($(0,3.3)!0.5!(0.5,6.5)$) {SPARQL};
  \fill[nbgreen!25] (0.8,6.6) rectangle (3,7.8);
  \node[font=\scriptsize] at ($(0.8,6.6)!0.5!(3,7.8)$) {Logic framework};
  \fill[nbteal!25] (1.4,7.9) rectangle (3,8.9);
  \node[font=\scriptsize] at ($(1.4,7.9)!0.5!(3,8.9)$) {Proof};
  \fill[nbpurple!20] (2.0,9.0) rectangle (3.9,10);
  \node[font=\scriptsize] at ($(2.0,9.0)!0.5!(3.9,10)$) {Trust};
  \fill[nbred!25] (3.1,3.3) rectangle (3.5,8.9);
  \node[font=\scriptsize, rotate=90] at ($(3.1,3.3)!0.5!(3.5,8.9)$) {Signature};
  \fill[nbred!25] (3.55,3.3) rectangle (3.95,8.9);
  \node[font=\scriptsize, rotate=90] at ($(3.55,3.3)!0.5!(3.95,8.9)$) {Encryption};
  \node[font=\small\bfseries] at (2,-0.8) {2005};
  \end{scope}
\end{tikzpicture}
\captionof{figure}{The Semantic Web stack according to Tim Berners-Lee, in 2003
(\emph{WWW past \& future}) and in 2005 (\emph{Web for real people}).}
\label{fig:sw-stack}
\end{center}

Reading the 2003 stack from the bottom:
\begin{itemize}[nosep]
  \item \textbf{URI} and \textbf{Unicode}: how to name any resource, how to write any
        character;
  \item \textbf{XML} and \textbf{namespaces}: a common syntax for structured documents,
        and a way to mix vocabularies without name clashes;
  \item \textbf{RDF} (model \& syntax) and \textbf{RDF Schema}: facts as triples, and a
        basic vocabulary of classes and properties
        (Section~\ref{sec:ontology-languages});
  \item \textbf{Ontology} (OWL): rich descriptions of the concepts of a domain;
  \item \textbf{Rules} and \textbf{logic framework}: inference beyond what the ontology
        language can express;
  \item \textbf{Proof} and \textbf{Trust}: explaining why a conclusion holds, and
        deciding whether to believe it (and its sources);
  \item \textbf{Signature} and \textbf{encryption} run alongside all the upper layers:
        who said what, and who can read it.
\end{itemize}

In the 2005 version the ontology layer is split into \textbf{OWL} and \textbf{Rules}
side by side, sharing a common fragment (\emph{DLP}, the part of OWL that can also be
written as rules), and the query language \textbf{SPARQL} appears next to them.

\begin{nbfact}{Stack or two towers?}
Horrocks, Parsia, Patel-Schneider and Hendler (\emph{Semantic Web Architecture: Stack or
Two Towers?}, 2005) argued that a single stack does not really work: rules with
\emph{negation as failure} (NAF, as in Prolog and Datalog) and the open-world logic of
OWL do not combine smoothly. They proposed \textbf{two towers} sharing the bottom layers
(URI, Unicode, XML, namespaces): one with \emph{DLP/Datalog}, rules and NAF, the other
with RDF, RDF Schema, OWL and first-order logic (FOL/SCL), possibly joined at the top by
a common logic framework. As an alternative they also drew a single stack in which the
rule layer (SWRL) sits \emph{above} OWL and below full FOL. The slides illustrate the
debate with Genoa landmarks: the two towers of Porta Soprana for the two-towers view,
the Lanterna (the lighthouse) for the single stack.
\end{nbfact}

\subsubsection{What is the Semantic Web really about?}
\slidesrc{Ontologies: 46}

James Hendler (2009) distinguishes three meanings in which the term is used:
\begin{itemize}[nosep]
  \item \textbf{Linked Data} (or Web of Data): datasets rendered into RDF, using very
        simple ontologies (really more like thesauri or vocabularies than AI languages)
        to provide terms for linking and describing the data;
  \item \textbf{Web 3.0}: sometimes used as a synonym for Semantic Web; it grew up as a
        way of explaining how semantics (especially the simple ontologies above) can be
        integrated with Web 2.0 technologies;
  \item \textbf{Semantic Web}: some people use it to mean all of the above (Hendler
        does).
\end{itemize}

\begin{nbwhy}
Linked Data is where the Semantic Web had its largest practical impact: resources such
as DBpedia and Wikidata are huge RDF knowledge graphs (Section~\ref{sec:what-is-an-ontology}),
and they are used in NLP for entity linking (connecting a name found by NER,
Section~\ref{sec:named-entity-recognition}, to the entity it denotes), question answering
and as background knowledge for language models.
\end{nbwhy}

\subsubsection{Web 4.0 and Web 5.0}
\slidesrc{Ontologies: 47--48}

\begin{itemize}
  \item \textbf{Web 4.0}: its services will be autonomous, proactive, self-learning,
        content-exploring, collaborative and content-generating \textbf{agents}, based on
        fully matured semantic and reasoning technologies as well as AI. They will
        support adaptive content presentation through an intelligent agent; examples are
        services interacting with sensors and implants, natural-language services and
        virtual reality services.
  \item \textbf{Web 5.0}: centred on users' \textbf{control over identity} (credentials
        for all platforms stored on a decentralised platform controlled only by the
        user, with no identity information collected by the Internet) and \textbf{data
        ownership} (complete decentralisation: users choose what to share with
        companies, and store their data with their own blockchain participation).
\end{itemize}
As the slides conclude, people no longer seem very interested in numbering the Web
versions.

\begin{nbsummary}[ontologies and the Semantic Web]
\begin{itemize}[nosep]
  \item \textbf{Ontology}: an explicit specification of a conceptualization (Gruber);
        classes, properties and relations with their meaning and constraints, at the
        ``semantic'' level, closer to first-order logic than database schemas. Main
        motivation: \textbf{semantic interoperability}.
  \item \textbf{Knowledge graph}: interlinked entities and facts described with an
        ontology; new facts can be inferred.
  \item \textbf{Methodology} (Noy \& McGuinness): scope, reuse, terms, classes and
        hierarchy, slots, facets, instances. Tool: \textbf{Protégé}.
  \item \textbf{RDF}: triples subject--predicate--object, forming a graph.
        \textbf{RDFS}: classes, \texttt{subClassOf}, properties with domain and range.
        \textbf{OWL}: richer primitives (disjointness, object vs.\ datatype properties,
        imports); Full (undecidable), DL (efficient reasoning), Lite (simplest).
  \item \textbf{Semantic Web}: the Web with machine-understandable meaning, organised as
        a layered stack (URI/XML $\to$ RDF $\to$ ontologies/rules $\to$ logic, proof,
        trust); in practice, mostly Linked Data.
\end{itemize}
\end{nbsummary}
